\chapter{Calibration}
\label{ch:calibracion}

Calibration is \textbf{the first step} of the analysis. It determines the
relationship between channel and velocity (mm/s) and, optionally, the isomer
shift reference. In Fitbauer the calibration is \textbf{fixed} and reused in all
subsequent loads, so it does not have to be redefined for each spectrum.

\section{Why calibration is needed}

The velocity transducer runs through a triangular cycle; the system records the
counts per channel. To convert the channel axis into a physical velocity axis
you need the \textbf{maximum velocity} \param{Vmax} (mm/s at the end of the
travel). This \param{Vmax} depends on the transducer amplitude and is common to
all spectra measured with the same configuration: this is why it is measured
once with a standard and reused.

The usual standard is \textbf{metallic $\alpha$-Fe}, whose sextet has very
well-known line positions. Fitbauer adopts the convention that an $\alpha$-Fe
spectrum corresponds to $\BHF = 33{,}0$~T, using the published velocity standard
($\pm0{,}839 / \pm3{,}084 / \pm5{,}329$~mm/s), just like NORMOS.

\begin{aviso}
The reference sextet positions are \textbf{not} derived from the nuclear
moments: the textbook calculation gives a splitting $\sim0{,}4\%$ smaller and
would make the $\BHF$ come out $\sim0{,}1$~T too high. The published $\alpha$-Fe
positions at 33.0~T are used instead.
\end{aviso}

\section{Defining a calibration from an $\alpha$-Fe}
\label{sec:definir-calib}

The usual procedure is to load an $\alpha$-Fe spectrum and use it to \textbf{fix
the velocity scale}. It is worth keeping one idea clear:

\begin{aviso}
\textbf{Calibrating means fitting the velocity, not the field.} In $\alpha$-Fe
the hyperfine field \textbf{is already known}: by convention $\BHF = 33{,}0$~T.
What is \emph{not} known is the \param{Vmax} —how many mm/s correspond to the end
of the transducer travel—. That is why the calibration fit \textbf{fixes}
$\BHF = 33{,}0$~T (and $\delta$, $\dEQ$ of $\alpha$-Fe) and \textbf{leaves the
\param{Vmax} free}, searching for the value that makes the six lines of the model
coincide with those of the spectrum. The result of the fit is not a $\BHF$ (that
is known), but the \param{Vmax} that defines the velocity axis for all subsequent
spectra.
\end{aviso}

\begin{enumerate}[nosep]
  \item \menu{File > Load…} and select the $\alpha$-Fe spectrum
        (for example \ruta{data\_sample/hierro\_metalico\_alphaFe.adt}).
  \item \textbf{Fix the standard and fit the velocity.} Set a single sextet with
        $\BHF = 33{,}0$~T and \textbf{mark it as fixed} (padlock), enable
        \emph{Fit Vmax} and run \menu{Fit > Fit}. The optimizer
        will move the \param{Vmax} (and refine $\delta$, $\Gamma$) until the six
        lines match. That \param{Vmax} value is the calibration.
  \item Mark the result as calibration with
        \menu{File > Use current file as calibration…}, or from the context
        menu of the file box (right click).
\end{enumerate}

\begin{nota}
If instead of fixing the field you leave $\BHF$ free, you will obtain a value
very close to 33~T (e.g.\ 33{,}04), but that is redundant: the standard
\emph{defines} 33{,}0~T. Fixing it avoids the degeneracy between $\BHF$ and
\param{Vmax} (both scale the line positions) and makes the fit determine
\emph{only} the velocity.
\end{nota}

\autofig{fig_calibration_fit}{Calibration of an $\alpha$-Fe: $\BHF = 33{,}0$~T
(standard) is \textbf{fixed} and the \param{Vmax} is \textbf{fitted}. The fit
returns $\text{Vmax} \approx 11{,}99$~mm/s with a reduced $\chi^2$ close to 1:
that \param{Vmax} is the calibrated velocity scale. Figure generated with
\ruta{scripts/make\_figures.py}.}

When marking the calibration, Fitbauer offers two variants from the context menu
of the file box:

\begin{itemize}[nosep]
  \item \menu{Use as calibration (current vmax and iso)}: takes directly the
        \param{Vmax} and the isomer shift $\delta$ of the current state.
  \item \menu{Use as calibration (details)…}: opens a dialog to review and
        edit the sample name, the \param{Vmax} and the IS before fixing.
\end{itemize}

\guicap{gui_use_as_calibration}{"Use as calibration" dialog with the sample name,
the calibrated Vmax and the reference isomer shift.}

\section{The calibration remains \emph{fixed}}
\label{sec:calib-fija}

Once defined, the calibration remains \textbf{fixed}: its \param{Vmax} (and the
reference IS) are saved and shown in the calibration label with the
\textbf{"fixed"} mark. From that point on it is applied automatically to the
spectra you load, according to this rule:

\begin{center}
\begin{tabular}{@{}p{0.46\textwidth}p{0.46\textwidth}@{}}
\toprule
\textbf{Action} & \textbf{Effect on the calibration} \\
\midrule
Load a file \textbf{without} its own calibration (raw counts:
\ruta{.ws5}, \ruta{.adt}) &
The \textbf{fixed one is used}: the velocity axis is built with the fixed
\param{Vmax}. The calibration \textbf{does not change}. \\[6pt]
Load a file \textbf{with} its own calibration (velocity axis:
\ruta{.csv}, \ruta{.txt}, \ruta{.dat}, \ruta{.exp}) &
It \textbf{replaces} the calibration with the file's own (its \param{Vmax}) and a
\textbf{warning} is shown in the status bar. \\
\bottomrule
\end{tabular}
\end{center}

That is: if you fix a calibration and then load \textbf{only data} of raw counts,
the calibration remains \textbf{intact} and is applied to all of them.

\begin{nota}
Raw-count files (\ruta{.ws5}/\ruta{.adt}) do not carry a velocity axis; they need
the calibration's \param{Vmax} to build it. Files already folded in velocity
(\ruta{.csv}, etc.) bring their own axis: that is why they \emph{replace} the
active calibration.
\end{nota}

\guicap{gui_calib_label_fixed}{Calibration label showing the "fixed" state, with
the sample, the Vmax and the reference IS.}

\subsection{Replacing the calibration}

When you load a file with its own calibration while another one is already fixed,
the new one replaces the previous one and Fitbauer shows a temporary warning in
the status bar of the type:

\begin{center}\ttfamily\small
Calibration replaced by the file's datos.csv: Vmax 12.0070 $\to$ 11.5000 mm/s
\end{center}

This is a record that the active \param{Vmax} has changed. The reference isomer
shift is preserved.

\subsection{Releasing the calibration}

To stop applying the fixed calibration, use \menu{Release fixed calibration} in
the context menu of the file box. From then on each load will use the
\param{Vmax} corresponding to its own file (or the control's default value).

\section{Calibrations from the web}

If you have access to the API, \menu{File > Web > Web calibrations…} lets you
download a remote calibration. Its values (\param{velocity\_calibrated} and
\param{isomer\_shift}) are fixed exactly like a local calibration, rescaling the
velocity axis of the current spectrum.

\section{Persistence in the session}

The fixed calibration is part of the \textbf{session}: when saving
(\menu{File > Save session…}) it is stored together with the rest of the state,
and when loading it (\menu{File > Load session…}) it is restored as is. In this
way an analysis is reproducible without recalibrating.

\begin{truco}
For batch work over many samples measured with the same configuration: load and
fix the $\alpha$-Fe calibration once, save it in a base session and start from it
each day.
\end{truco}
